\documentclass[11pt]{article}

\usepackage{amsmath,amssymb,amsthm}
\usepackage[margin=1in]{geometry}
\usepackage{tikz}
\usetikzlibrary{positioning}

\newtheorem{theorem}{Theorem}
\newtheorem{lemma}{Lemma}
\newtheorem{corollary}{Corollary}
\theoremstyle{definition}
\newtheorem{definition}{Definition}
\theoremstyle{remark}
\newtheorem{remark}{Remark}

\newcommand{\Dcal}{\mathcal{D}}
\newcommand{\Scal}{\mathcal{S}}
\newcommand{\Acal}{\mathcal{A}}
\newcommand{\Zcal}{\mathcal{Z}}

\title{Every Good Regulator of a System Must Be a Model of That System:\\
\large An Information-Theoretic Reformulation}
\author{Neil D. Lawrence}
\date{\today}

\begin{document}
\maketitle

\begin{abstract}
Conant and Ashby's (1970) good regulator theorem is remembered mainly for its
slogan. The paper itself buries the argument under Sommerhoff's five-set
formalism and a long detour through group homomorphisms, machine
isomorphisms, and ``Black Box'' equivalence. Yet the proof of the theorem
proper is already, in essence, an entropy argument: Conant and Ashby adopt
$H(Z)$, the entropy of the outcome, as their criterion of success, and the
proof turns entirely on how the entropy of a mixture responds to
concentrating probability mass. This note rewrites the paper end to end in
the language of probability and information theory that would be familiar
to anyone who has read MacKay (2003) or Cover and Thomas (2006): random
variables in place of Sommerhoff's sets, conditional distributions --
policies, in the language of reinforcement learning -- in place of Ashby's
relations, and entropy, conditional entropy and mutual information in place
of homomorphisms and isomorphisms. What falls out is a shorter derivation
and a sharper statement of the theorem: an optimal regulator's action must
have \emph{zero entropy conditioned on the state of the system it
regulates}. That single condition is what it means, precisely and only,
for a regulator to be a model of its system.
\end{abstract}

\section{Introduction}

Regulation is the business of keeping some outcome close to a goal in the
face of disturbance. Conant and Ashby's question was whether a regulator
that succeeds at this must, of necessity, contain something that deserves
to be called a model of the system it regulates -- not as a design choice,
but as a logical consequence of being good at the job.

Their answer was yes, and the proof is short once the right criterion of
success is fixed. They fix it early: the entropy $H(Z)$ of the outcome
variable is to be made as small as possible. Everything that follows in
their paper, once this is granted, is a fact about how entropy behaves
under mixtures -- except for a lengthy central section in which they try,
and by their own admission fail, to pin down what ``model'' should mean by
appeal to isomorphisms of groups and automata. That section is the one
piece of the paper that is not really about information at all, and it is
the piece we replace here.

The version below keeps every substantive claim of the original but states
the whole thing -- setting, success criterion, the error/cause-controlled
distinction, the notion of a model, and the theorem's proof -- in terms of
random variables and their entropies. We assume the reader already knows
what entropy, conditional entropy, and mutual information are, together
with their standard properties (the chain rule, non-negativity, and the
data-processing inequality); see MacKay (2003) or Cover and Thomas (2006)
for the definitions if a refresher is wanted. We leave the base of the
logarithm unspecified, as Conant and Ashby do: nothing below depends on
whether entropy is measured in bits or nats. One further notational
choice: we write the regulator's action as $A$ rather than Ashby's $R$,
both to match the state/action pairing familiar from reinforcement
learning ($s$, $a$, $\pi(a\mid s)$) and to leave the symbol $r$ free for a
reward function introduced in Section~\ref{sec:setup}.

\section{Regulation as a Probabilistic System}
\label{sec:setup}

Sommerhoff's apparatus of five interacting sets $Z, G, R, S, D$ translates
directly into four random variables and a goal set:

\begin{itemize}
\item $D$, taking values in a finite alphabet $\Dcal$: the disturbance --
  the exogenous, unpredictable factor (weather, demand, the bird's flight
  path) that regulation exists to counteract.
\item $S$, taking values in $\Scal$: the state of the \emph{reguland}, the
  part of the system other than the regulator (aircraft positions, fuel
  levels; the dynamics of gun and bird) -- in reinforcement-learning terms,
  simply the state of the environment.
\item $A$, taking values in $\Acal$: the regulator's action (settings in
  the control tower; the shot the hunter fires) -- the action, in
  reinforcement-learning terms.
\item $Z$, taking values in $\Zcal$: the outcome, determined jointly by $S$
  and $A$ through a (possibly stochastic) channel $p(z\mid s,a)$. Where
  convenient we take this channel to be a deterministic function
  $z=\psi(s,a)$, as Conant and Ashby do; nothing below hinges on this,
  since the argument goes through with $p(z\mid s,a)$ in place of
  $\mathbb{1}[z=\psi(s,a)]$ throughout.
\item $G\subset\Zcal$: the set of ``good'' outcomes.
\end{itemize}

A joint distribution $p(d,s,a,z)$ is assumed to exist, factorising
according to the causal structure of the problem (Section~\ref{sec:causal}
makes this concrete). Regulation is the choice of a conditional
distribution over actions given state -- a \emph{policy}
\[
\pi(a\mid s) := p(a\mid s),
\]
to use the term the way reinforcement learning does (Sutton and Barto,
1998) -- with $p(d)$ and hence $p(s)$ taken as given and fixed. Two
features of this setting are not standard reinforcement learning, and are
worth flagging before the analogy is pushed any further. First, there is
no time index and no dynamics: this is a single decision, not a
trajectory, so nothing here resembles a value function, a Bellman
equation, or discounting. Second, and more importantly, there is not yet a
reward -- only an outcome $Z$ whose entropy we are about to ask to be
small. We return to the relationship between the two immediately below,
and again in Section~\ref{sec:theorem}.

\paragraph{From a goal set to an entropy.} Sommerhoff's own criterion of
success is set membership: regulation succeeds when $Z\in G$ with
probability~1. This is also the natural point at which to define a
reward, in the reinforcement-learning sense:
\[
r(z) := \mathbb{1}[z\in G], \qquad\text{so that}\qquad \mathbb{E}[r(Z)] = \Pr(Z\in G),
\]
and the standard reinforcement-learning objective would be to choose
$\pi(a\mid s)$ to maximise $\mathbb{E}[r(Z)]$. A weaker, numerical version
of Sommerhoff's criterion asks instead that some scalar loss have small
root-mean-square. Conant and Ashby's own choice -- and ours -- is neither
of these but entropy: regulation succeeds to the extent that $H(Z)$ is
small, and \emph{optimal} regulation is regulation that minimises $H(Z)$
given $p(s)$ and the channel $p(z\mid s,a)$.

This is a different, and in general strictly stronger, criterion than
reward maximisation. Maximising $\mathbb{E}[r(Z)]$ only asks that $Z$ land
in $G$ with high probability; it says nothing about which outcome within
$G$, or within its complement, results, nor how reliably. Minimising
$H(Z)$ asks that the whole distribution of $Z$ be as concentrated as the
channel $p(z\mid s,a)$ allows, wherever it is concentrated. The two
criteria coincide exactly when some policy can place all its mass on a
single outcome inside $G$; whenever the achievable floor of $H(Z)$ is
strictly positive, reward maximisation is silent about how that residual
uncertainty is distributed, while the entropy criterion pins it down
completely. We return to this gap, and to why we do not close it by
exponentiating $r$, in Section~\ref{sec:theorem}.

The entropy criterion's advantage over both alternatives is that it is
defined whenever $\Zcal$ is a set of classes, with no need for $\Zcal$ to
carry a metric or an additive structure -- Conant and Ashby's examples are
species of fish caught in trawling and amino-acid chains produced by a
ribosome, neither of which supports a root-mean-square.

We write
\[
H(Z) := -\sum_{z\in\Zcal} p(z)\log p(z),
\]
and use $H(\cdot\mid\cdot)$ and $I(\cdot\,;\cdot)$ for conditional entropy
and mutual information throughout.

\section{Two Channels From Disturbance to Regulator}
\label{sec:causal}

Conant and Ashby distinguish \emph{cause-controlled} regulation, in which
the regulator responds directly to the disturbance $D$, from
\emph{error-controlled} regulation, in which it responds only to the
state $S$ that the disturbance has already partly corrupted. The
distinction is a statement about which Markov chain the joint distribution
respects, and the consequence is a direct application of the data
processing inequality.

\paragraph{Cause-controlled.} Here $S$ and $A$ are separate, possibly
stochastic, functions of the common cause $D$:
\[
p(d,s,a,z) = p(d)\,p(s\mid d)\,p(a\mid d)\,p(z\mid s,a).
\]

\begin{center}
\begin{tikzpicture}[node distance=1.6cm, >=stealth]
\node (D) {$D$};
\node (S) [above right=8mm and 18mm of D] {$S$};
\node (A) [below right=8mm and 18mm of D] {$A$};
\node (Z) [right=32mm of D] {$Z$};
\draw[->] (D) -- (S);
\draw[->] (D) -- (A);
\draw[->] (S) -- (Z);
\draw[->] (A) -- (Z);
\end{tikzpicture}
\end{center}

Here $A$ has direct access to $D$, so $I(D;A)$ can in principle be as large
as $H(D)$: if $A$ is a sufficient statistic of $D$ for $Z$'s purposes,
$H(Z)$ can be driven all the way to its floor, $H(Z\mid D)$ -- in the
favourable case where $\psi$ makes that floor zero, regulation can be made
perfect.

\paragraph{Error-controlled.} Here $A$ has no direct channel to $D$; it
sees only $S$, which has already been shaped by $D$:
\[
p(d,s,a,z) = p(d)\,p(s\mid d)\,p(a\mid s)\,p(z\mid s,a),
\]
i.e.\ $D\to S\to A$ is a Markov chain.

\begin{center}
\begin{tikzpicture}[node distance=1.6cm, >=stealth]
\node (D) {$D$};
\node (S) [right=22mm of D] {$S$};
\node (A) [below=10mm of S] {$A$};
\node (Z) [right=22mm of S] {$Z$};
\draw[->] (D) -- (S);
\draw[->] (S) -- (A);
\draw[->] (S) -- (Z);
\draw[->] (A) -- (Z);
\end{tikzpicture}
\end{center}

By the data processing inequality applied to $D\to S\to A$,
\begin{equation}
I(D;A) \;\le\; I(D;S).
\label{eq:dpi}
\end{equation}
The regulator can never know more about the disturbance than the state $S$
already reveals -- and $S$, being the very variable whose fluctuation
regulation is meant to suppress, has typically already let some of that
disturbance through into $Z$ by the time $A$ acts on it. Conant and
Ashby's example is the cow: an error-controlled thermoregulatory reflex
only fires once blood temperature has already fallen, because the
reflex's channel to the cause (cold air) runs entirely through the effect
(temperature) it is trying to prevent. A cause-controlled system -- a
temperature sensor on the animal's skin, reacting to the cold air stream
itself -- can in principle correct before any temperature drop occurs at
all. Inequality \eqref{eq:dpi} is the general reason error control is, as
Conant and Ashby put it, ``a primitive and demonstrably inferior method of
regulation'': whatever the residual $H(Z\mid D)$ available to cause
control, error control can only ever access an amount of information about
$D$ bounded by what already leaked into $S$, so its achievable $H(Z)$ is
bounded away from that floor in general.

From here on, following Conant and Ashby, we consider regulation of the
more general kind in which $\pi(a\mid s)$ is free to be chosen without
restriction -- this covers cause control as a special case (via $D\to S$
and $D\to A$ jointly determining an effective $\pi(a\mid s)$) and includes
error control as the case actually analysed in the theorem below.

\section{What Is a Model? An Entropy Account}
\label{sec:models}

Conant and Ashby spend the middle third of their paper trying to fix a
notion of ``model'' by appeal to isomorphism, working through finite-group
homomorphisms, Hartmanis and Stearns' machine homomorphisms, and Black-Box
equivalence, only to conclude that there are about as many possible
definitions of isomorphism as there are of model, with no natural boundary
between the two. Their retreat is to define a model \emph{operationally},
via the mapping $h$ that falls out of the theorem itself.

We can do better than retreat, because entropy gives a non-arbitrary family
of definitions with a built-in ordering, and the theorem (Section
\ref{sec:theorem}) will then tell us exactly which member of the family is
forced by optimality.

\begin{definition}[Model]
$A$ is a \emph{model} of $S$ if $A$ is a deterministic function of $S$:
\[
H(A\mid S) = 0.
\]
\end{definition}

$H(A\mid S)=0$ is precisely the statement that there exists a function
$h:\Scal\to\Acal$ with $A=h(S)$ almost surely -- so this recovers Conant
and Ashby's ``mapping $h:S\to R$'' language, but as a \emph{consequence} of
a condition on entropy rather than as a primitive assumption borrowed from
group theory.

\begin{definition}[Faithful model]
$A$ is a \emph{faithful} (or \emph{lossless}) model of $S$ if
\[
H(A\mid S) = H(S\mid A) = 0,
\]
equivalently $I(A;S) = H(A) = H(S)$.
\end{definition}

A faithful model is invertible: $S$ and $A$ carry exactly the same
information, related by a bijection $h$. This is the information-theoretic
counterpart of the group-theoretic isomorphism Conant and Ashby start from
in Section~4 of the original -- but stated without any need for $S$ or $A$
to carry group structure at all.

\begin{definition}[Lossy model]
\label{def:lossy}
$A$ is a \emph{lossy} (or \emph{homomorphic}) model of $S$ if
\[
H(A\mid S) = 0 \quad\text{but}\quad H(S\mid A) > 0.
\]
\end{definition}

Here $h$ is many-to-one: several states $s$ collapse onto the same action
$a$. This is the counterpart of Conant and Ashby's machine homomorphism --
a structure-preserving map that discards some detail -- again recovered
without any appeal to algebraic structure, purely from the size of
$H(S\mid A)$.

\begin{remark}
The three definitions above sit inside a single one-parameter family
indexed by $H(A\mid S)\ge 0$, running from a model at one end
($H(A\mid S)=0$) through arbitrarily loose statistical association at the
other ($H(A\mid S)$ close to $H(A)$, i.e.\ $A$ nearly independent of $S$).
This is exactly the continuum Conant and Ashby worried about when they
observed that homomorphism definitions could be weakened indefinitely with
``no natural boundary dividing model from non-model.'' Fixing $H(A\mid
S)=0$ as \emph{the} definition of model looks, on its own, like one more
arbitrary point on that continuum. What removes the arbitrariness is the
theorem of Section~\ref{sec:theorem}: optimality in the sense of minimal
$H(Z)$ forces $H(A\mid S)=0$ exactly, not approximately and not as a matter
of degree. The boundary Conant and Ashby went looking for among
isomorphism definitions was never going to be found there; it is a
consequence of the regulation problem, not a property of any particular
mapping.
\end{remark}

\begin{remark}[Connection to sufficient statistics]
A lossy model $A=h(S)$ that nonetheless preserves everything about $S$
relevant to some downstream variable $Z$ -- i.e.\ $I(A;Z)=I(S;Z)$ while
$H(A\mid S)=0$ -- is precisely a \emph{sufficient statistic} of $S$ for
$Z$, in Fisher's classical sense. Seeking the coarsest (lowest-entropy)
such $A$ is the same problem addressed, in a Lagrangian relaxation, by the
information bottleneck method (Tishby, Pereira and Bialek, 1999): minimise
$I(A;S)$ subject to a constraint on $I(A;Z)$. The good regulator theorem
below is the limiting case of that trade-off in which the relevance
constraint is non-negotiable ($H(Z)$ must be at its minimum) and only
simplicity is optimised subject to it.
\end{remark}

\section{The Good Regulator Theorem}
\label{sec:theorem}

We now state and prove the theorem in entropy form. The setting is that of
Section~\ref{sec:setup}: finite alphabets $\Scal,\Acal,\Zcal$, a fixed
distribution $p(s)$, a fixed (for now, deterministic) outcome map
$\psi:\Scal\times\Acal\to\Zcal$, and a free choice of policy $\pi(a\mid
s)$. Write
\[
\Pi := \operatorname*{arg\,min}_{\pi(a\mid s)} H(Z)
\]
for the class of optimal policies. As in the original, we set aside the
possibility of several distinct optimal marginals $p(z)$ tying for the
minimum $H(Z)$; this is a genuine possibility but complicates the proof
without changing its content, so we assume every member of $\Pi$ induces
the same $p(z)$.

\begin{lemma}[Concentration lemma]
\label{lem:concentration}
For every $\pi(a\mid s)\in\Pi$ and every $s$ with $p(s)>0$,
\[
H(Z\mid S=s) = 0.
\]
Equivalently: for every optimal policy, all actions $a$ with
$\pi(a\mid s)>0$ satisfy $\psi(s,a)=z^\ast(s)$ for a single outcome
$z^\ast(s)$ depending on $s$ alone.
\end{lemma}

\begin{proof}
Suppose, for contradiction, that some $\pi(a\mid s)\in\Pi$ and some $s$
with $p(s)>0$ admit $a_1\ne a_2$ with $\pi(a_1\mid s)>0$, $\pi(a_2\mid
s)>0$, and $\psi(s,a_1)=z_1\ne z_2=\psi(s,a_2)$.

Consider transferring probability $\alpha$ from the pair $(s,a_1)$ to the
pair $(s,a_2)$, holding $\pi(a\mid s')$ fixed for every $s'\ne s$ and every
other value of $a$ at $s$ fixed. This changes only the marginal
probabilities of $z_1$ and $z_2$, replacing $(p(z_1),p(z_2))=(u,v)$ by
$(u-\alpha,\,v+\alpha)$, for $\alpha$ ranging over an interval containing
$0$ in its interior (since $\pi(a_1\mid s),\pi(a_2\mid s)>0$). All other
terms of $H(Z)$ are unaffected. The two affected terms are
\[
g(\alpha) = -(u-\alpha)\log(u-\alpha) - (v+\alpha)\log(v+\alpha),
\]
a strictly concave function of $\alpha$
($g''(\alpha) = -\tfrac{1}{u-\alpha}-\tfrac{1}{v+\alpha} < 0$), maximised
uniquely at $\alpha^\ast = (u-v)/2$. Since $\alpha=0$ lies in the interior
of the admissible interval, either $\alpha=0=\alpha^\ast$, in which case
\emph{any} admissible nonzero $\alpha$ strictly decreases $g$, or
$\alpha=0\ne\alpha^\ast$, in which case moving $\alpha$ further away from
$\alpha^\ast$ (which remains admissible for small enough moves) strictly
decreases $g$. Either way there is an admissible perturbation that
strictly decreases $H(Z)$, contradicting optimality of the original
$\pi(a\mid s)$. Hence no such $a_1,a_2$ exist, which is exactly the claim.
\end{proof}

\begin{theorem}[Good regulator theorem, entropy form]
\label{thm:goodreg}
There exists an optimal policy $\pi(a\mid s)\in\Pi$ that is a
\emph{model} of $S$ in the sense of Section~\ref{sec:models}:
\[
H(A\mid S) = 0.
\]
Equivalently, there is a function $h:\Scal\to\Acal$ such that setting
$A=h(S)$ achieves the minimal outcome entropy $H(Z)$.
\end{theorem}

\begin{proof}
Take any $\pi(a\mid s)\in\Pi$. By Lemma~\ref{lem:concentration}, for each
$s$ with $p(s)>0$ every $a$ with $\pi(a\mid s)>0$ maps, together with $s$,
to the same outcome $z^\ast(s)$. Choose, for each such $s$, one particular
such $a$ and call it $h(s)$. Define $\pi'(a\mid s) := \mathbb{1}[a=h(s)]$.

Because every $a$ with $\pi(a\mid s)>0$ under the original policy already
produced $z^\ast(s)$, redistributing all of that mass onto the single
value $h(s)$ leaves $p(z^\ast(s))$, and hence the entire marginal $p(z)$,
unchanged. So $\pi'(a\mid s)$ induces the same $H(Z)$ as the policy we
started from, and therefore $\pi'(a\mid s)\in\Pi$. By construction
$H'(A\mid S)=0$, since $A$ is exactly the deterministic function $h$ of
$S$ -- a deterministic policy, in reinforcement-learning terms.
\end{proof}

\begin{corollary}
The theorem's regulator is, in general, a \emph{lossy} model of $S$ rather
than a faithful one: any two states $s,s'$ with $z^\ast(s)=z^\ast(s')$ that
admit a common suitable response can be assigned the same $h(s)=h(s')$
without affecting $H(Z)$. The coarsest such $h$ -- merging every pair of
states that require the same regulatory response -- is the minimal
sufficient statistic of $S$ for the regulation task, in the sense of the
remark closing Section~\ref{sec:models}.
\end{corollary}

\subsection*{Remarks}

\begin{itemize}
\item \textbf{Unnecessary complexity has a name.} Lemma~\ref{lem:concentration}
holds for \emph{every} member of $\Pi$, not only the deterministic one
constructed in the proof: every optimal policy already has
$H(Z\mid S)=0$. What it need not already have is $H(A\mid S)=0$: a policy
can be optimal while carrying extra randomness in $A$ that varies without
ever changing which $z^\ast(s)$ results. That surplus, $H(A\mid S)$
itself, is a direct entropy measure of the ``unnecessary complexity''
Conant and Ashby gesture at informally. Not every optimal regulator is a
model of its reguland; the ones that are not are exactly the ones
carrying positive $H(A\mid S)$ that the task does not require.

\item \textbf{The search space collapses.} Once Theorem~\ref{thm:goodreg}
is granted, searching for a good regulator is a search over deterministic
functions $h:\Scal\to\Acal$, not over conditional distributions
$\pi(a\mid s)$ -- a search over $|\Acal|^{|\Scal|}$ candidates rather than
a continuum.

\item \textbf{Cause- versus error-controlled regulation revisited.} In the
error-controlled setting of Section~\ref{sec:causal}, $A$'s only input is
$S$, and the theorem applies directly: $A=h(S)$. In the cause-controlled
setting, $S$ and $A$ are both functions of the shared cause $D$
($S=\phi(D)$, $A=\rho(D)$ in the deterministic case), so the theorem's
$h$ must additionally satisfy $h(\phi(d))=\rho(d)$ for all $d$ -- $h$ and
$\rho$ must \emph{agree through} $\phi$. This is a strictly stronger
demand: not only must $A$ be recoverable from $S$, but the particular
recovering function $h$ must coincide with $A$'s own direct route from
$D$. It is the information-theoretic reading of Conant and Ashby's remark
that Fig.~1's regulator ``must be a homo- or isomorph of $S$,'' where
Fig.~2's need only be $A=h(S)$.

\item \textbf{Time-varying regulands.} If $p(s)$ drifts slowly, the
argument holds over any window in which $p(s)$ is effectively constant,
with $h$ tracking the drift. The regulator is then an adaptive, rather
than a fixed, model of the reguland -- exactly the situation of online
estimation or filtering, where a sufficient statistic must itself be
updated as the underlying distribution changes.
\end{itemize}

\subsection*{Reward, inference, and control}

Section~\ref{sec:setup} defined a reward $r(z):=\mathbb{1}[z\in G]$ purely
so that the standard reinforcement-learning objective,
$\max_\pi\mathbb{E}[r(Z)]=\max_\pi\Pr(Z\in G)$, would be available for
comparison. It is not what Theorem~\ref{thm:goodreg} optimises, and the
gap between the two objectives is informative rather than incidental.
Maximising $\mathbb{E}[r(Z)]$ only asks that $Z$ land in $G$ with high
probability; it is silent about which outcome within $G$ (or within its
complement) results, and how reliably. Minimising $H(Z)$ asks for the
stronger thing established by Lemma~\ref{lem:concentration}: that wherever
the outcome distribution sits, it be as concentrated as the channel
$p(z\mid s,a)$ allows. The two objectives agree exactly when some policy
can put all its probability mass on a single outcome in $G$; whenever the
achievable floor of $H(Z)$ is strictly positive, reward maximisation has
nothing further to say about how that residual uncertainty is shaped,
while the entropy criterion pins it down completely.

This is close in spirit to, but should not be confused with, the
\emph{control as inference} framework (Todorov, 2008; Toussaint, 2009;
Rawlik, Toussaint and Vijayakumar, 2013; Levine, 2018, gives a careful
tutorial). That framework also turns control into an inference problem,
but does so by introducing an auxiliary binary ``optimality'' variable $O$
with an exponentiated-reward observation model,
$p(O=1\mid s,a)\propto\exp\!\big(r(s,a)/\alpha\big)$, and then computing
$p(a\mid s, O=1)$. The construction works, but it has known costs: the
exponential transform and the temperature $\alpha$ are modelling choices
with no derivation behind them, and under stochastic dynamics the
resulting policy can be systematically optimistic, effectively taking
credit for good luck in transitions it does not control, unless the
inference is carried out carefully (Levine, 2018, \S5).

None of that machinery is needed here, and it is worth being explicit
about why. The good regulator theorem never introduces an auxiliary
optimality variable to condition on: its ``evidence'' is $Z$ itself, an
outcome that actually occurs, and its objective is the ordinary Shannon
entropy of that outcome, not a KL divergence to an exponentiated-reward
posterior. There is no reward to exponentiate and no temperature to fix.

If an analogy is wanted at all, it is not to another inference framework
but to a representational one: O'Regan's sensorimotor contingency theory
(O'Regan, 1992; O'Regan and No\"e, 2001). That work argues, on entirely
independent grounds, that perceiving does not require the nervous system
to build and maintain a detailed internal replica of the scene; it
requires mastery of the lawful contingencies linking action to the
sensory consequences of action, with whatever the mapping does not need
left in the world to be resampled through the loop rather than stored --
O'Regan's phrase is that the world serves as ``an outside memory.''
Theorem~\ref{thm:goodreg} and its Corollary say the same thing in entropy
terms. The good regulator's action is a deterministic function of the
sensed state, $A=h(S)$, but $h$ is generically many-to-one, so the
regulator carries only the minimal sufficient statistic of $S$ needed to
act, not a faithful copy of $S$ and certainly not of the disturbance $D$
that produced it. What the theorem calls a lossy model
(Definition~\ref{def:lossy}) is what O'Regan's account calls mastery of a
sensorimotor contingency, and the two identifications need no
exponentiated reward, no temperature, and no auxiliary optimality
variable: the causal loop of Section~\ref{sec:causal},
$D\to S\to A\to Z$, already gives $S$ and $A$ the sensory-state and action
roles the analogy needs. The chronology also runs the opposite way from
the usual story about this kind of idea: Conant and Ashby's theorem
predates O'Regan's argument for the world as an outside memory by over
twenty years, and O'Regan's argument in turn predates control as
inference and the free-energy literature by well over a decade -- none of
which this reading needs to borrow from.

\section{Discussion}

Section~\ref{sec:models} left a genuine puzzle unresolved: entropy gives a
whole family of candidate definitions of ``model,'' ordered by
$H(A\mid S)$, and nothing internal to that family picks out one point as
special. Theorem~\ref{thm:goodreg} is what picks it out. It is not that
$H(A\mid S)=0$ is a reasonable place to draw the line; it is that
minimising $H(Z)$ \emph{forces} $H(A\mid S)$ to zero, exactly, for at least
one member of every optimal class. Where Conant and Ashby's original
Section~4 left ``model'' underdetermined by algebra, the theorem
overdetermines it by information: a good regulator's action, in the
optimal case, does not merely resemble the reguland's state -- it is a
deterministic (generally lossy) readout of it, in precisely the sense that
$A$ carries zero residual uncertainty once $S$ is known.

This changes model-building from an optional design choice to a logical
consequence of optimal regulation, in exactly the sense Conant and Ashby
intended: a controller, a Kalman filter, or a reinforcement-learning
policy $\pi(a\mid s)$ that succeeds at minimising the entropy of some
outcome variable it does not fully control has, whether its designer
intended it or not, come to encode a compressed, deterministic function of
the state of the system it is acting on. It has done so for
information-theoretic reasons, not because anyone told it to.

For the study of the brain, this reading of the theorem sharpens rather
than only repeats Conant and Ashby's own conclusion. The claim is not the
qualitative one that the brain probably models its environment, nor even
that it must model its environment in some loose isomorphism-flavoured
sense; it is the quantitative one that whatever mapping from sensory state
to motor or physiological action the nervous system settles on, to the
extent that mapping is entropy-optimal it will exhibit $H(A\mid S)=0$: no
randomness in the response beyond what the sensed state already
determines. That is an empirically meaningful claim about neural
variability, not only a metaphor, and it sits comfortably alongside later,
independently motivated proposals -- efficient coding (Barlow, 1961), the
information bottleneck (Tishby, Pereira and Bialek, 1999), and
sensorimotor contingency theory's argument that perception need not
reconstruct the scene at all, only the action-contingent statistics of it
(O'Regan and No\"e, 2001) -- that arrive at similar demands for
compressed, task-sufficient internal representations from different
starting points.

\section*{References}

\begin{itemize}
\item Barlow, H.~B. (1961). Possible principles underlying the
  transformations of sensory messages. In \emph{Sensory Communication},
  ed.\ W.~Rosenblith, MIT Press.
\item Conant, R.~C. and Ashby, W.~R. (1970). Every good regulator of a
  system must be a model of that system. \emph{International Journal of
  Systems Science}, 1(2):89--97.
\item Cover, T.~M. and Thomas, J.~A. (2006). \emph{Elements of Information
  Theory}, 2nd edition. Wiley-Interscience.
\item Levine, S. (2018). Reinforcement learning and control as
  probabilistic inference: tutorial and review. \emph{arXiv:1805.00909}.
\item MacKay, D.~J.~C. (2003). \emph{Information Theory, Inference, and
  Learning Algorithms}. Cambridge University Press.
\item O'Regan, J.~K. (1992). Solving the ``real'' mystery of visual
  perception: the world as an outside memory. \emph{Canadian Journal of
  Psychology}, 46(3):461--488.
\item O'Regan, J.~K. and No\"e, A. (2001). A sensorimotor account of
  vision and visual consciousness. \emph{Behavioral and Brain Sciences},
  24(5):939--973.
\item Rawlik, K., Toussaint, M. and Vijayakumar, S. (2013). On stochastic
  optimal control and reinforcement learning by approximate inference. In
  \emph{Proceedings of Robotics: Science and Systems (RSS)}.
\item Sommerhoff, G. (1950). \emph{Analytical Biology}. Oxford University
  Press.
\item Sutton, R.~S. and Barto, A.~G. (1998). \emph{Reinforcement Learning:
  An Introduction}. MIT Press.
\item Tishby, N., Pereira, F.~C. and Bialek, W. (1999). The information
  bottleneck method. In \emph{Proceedings of the 37th Annual Allerton
  Conference on Communication, Control, and Computing}, pp.~368--377.
\item Todorov, E. (2008). General duality between optimal control and
  estimation. In \emph{Proceedings of the 47th IEEE Conference on
  Decision and Control (CDC)}.
\item Toussaint, M. (2009). Robot trajectory optimization using
  approximate inference. In \emph{Proceedings of the 26th International
  Conference on Machine Learning (ICML)}.
\end{itemize}

\end{document}
